The {\bf Net} class is small, its main function being to hold a net
identifier and associated NET class instance. The net connections to
elements are instances of class {\bf Net} (a net identifier wrapper),
not class {\bf NET} (the actual net). When two {\bf Net}s are connected
together it is not necessary to visit all elements to which the {\bf
Net}s are connected, instead one of the associated {\bf NET} class
instances is merged into the other and both {\bf Net} instances now
contain the same {\bf NET} instance.

A {\bf Net} constructor which is not passed an identifier will create
one consisting of "n" with an appended unique integer.

The member variables for {\bf Net} are -

{\bf private String id} - the net identifier.

{\bf private NET net} - the actual net.

There are six other member variables, input\_elements, input\_ports, output\_elements,
output\_ports, input\_index and output\_index, which are just temporary variables
used when traversing connected elements.

There are some class variables -

{\bf private static int icount} - an integer used to form unique
net identifiers of the form "n123". It is incremented after each use.

{\bf private static int hashcode\_count} - an integer used to form a
crude unique hash code for use by methods {\bf equals()} and {\bf
hashCode()} on instances of class {\bf NET}. These methods are used by
HashMap and HashSet classes to distinguish or equate NET class
instances.

{\bf private static HashMap\verb'<'String, Net\verb'>' identmap = new HashMap\verb'<'String,
Net\verb'>'(20000)} - a map of nets. The key is a net identifier and the
associated value is class Net.

{\bf private static HashSet\verb'<'NET\verb'>' instances = new HashSet\verb'<'NET\verb'>'(20000)} -
the set of NETs. This set is traversed when writing nets to the netlist
file.

{\bf public static Net LO} \\
{\bf public static Net HI} -
globally available GND and VCC nets. {\bf LO} and {\bf HI} are assigned
in the constructor XElements() because they are then connected there to
newly constructed Elements GND and VCC.




For an operation on a net the net identifier is used to extract the
{\bf Net} class instance from static map {\bf instances}. The {\bf Net}
in turn contains variable {\bf net}, class {\bf NET}, which is the
actual net. Many instances of net identifier class {\bf Net} will
contain the same actual {\bf NET}, meaning that they are connected.

The main methods are described in the following subsections.

\subsection{constructors and class creation methods}

    Constructor {\bf Net()} creates an un-named net with a unique
    identifier of the form "n123".

    A net with a constructed identifier can also be produced by
    constructor
    \begin{lstlisting}[gobble=4, language=java]
        public Net (String header)
    \end{lstlisting}
    which creates a net with a unique identifier by concatenating a
    unique integer to the header string supplied. This is used for
    signals which are not derived from source code variables but which
    are easier to trace in the netlist if given meaningfull identifier
    headers.

    A named Net can be constructed using constructor
    \begin{lstlisting}[gobble=4, language=java]
        public Net (String ident, boolean large)
    \end{lstlisting}
    which creates a net using the supplied identifier. The second
    parameter {\bf large}, if true, causes lists inside class
    NET to be created with a large initial size (10000) instead of the
    default size. This is used for Nets {\bf GND} and {\bf VCC} as they have
    a very large number of connections.

    A named net can also be constructed using method
    \begin{lstlisting}[gobble=4, language=java]
        static public Net get (String ident)
    \end{lstlisting}
    which returns a named net if it already exists and otherwise
    constructs it.

    An array of unnamed Nets can be created by method
    \begin{lstlisting}[gobble=4, language=java]
        static public Net[] netArray (int size)
    \end{lstlisting}
    where parameter {\bf size} provides the dimension of the array.

\subsection{method init()}

    This method re-initialises the class static variables prior
    to any run of 3PL.

\subsection{method equals()}

    \begin{lstlisting}[gobble=4, language=java]
        public boolean equals (Object o)
    \end{lstlisting}
    This method returns true if the argument is the same Net or
    is connected to the object Net (they share the same NET class
    instance). If the argument is null or not class Net then
    false is returned.

\subsection{connection to an Element}

    \begin{lstlisting}[gobble=4, language=java]
        public void connect (Element e, String port, PortType type)
    \end{lstlisting}
    This method connects a Net to an Element port but only modifies the
    Net structure, not that of the Element to which it connects. This is
    only called from methods {\bf addInput()}, {\bf addOutput()} or {\bf
    addTS()} in class Element which handle the related changes in that
    class.

\subsection{disconnection from an Element}

    \begin{lstlisting}[gobble=4, language=java]
        public void disconnect (Element e, String port)
    \end{lstlisting}
    This method disconnects a Net to an Element port but only modifies
    the Net structure, not that of the Element to which it is connected.
    The method is only called in -
    \blist
    \item   Element.strip()
    \item   Element.removeNet()
    \item   TDECode.optimiseXOR()
    \item   TDECode.inputEliminate()
    \blend
    In each case related changes to the Element are carried out
    elsewhere


\subsection{connection to another Net}

    \begin{lstlisting}[gobble=4, language=java]
        public void connect(Net n) {
            if (n.equals(LO) || n.equals(HI)) {
                n.connect(this);
                return;
            }

            if (this.equals(n))
                throw new ExEx("SYSTEM ERROR - netlist self-connect");

            if ((getNumOutputs() + n.getNumOutputs()) > 1) {
                msg("connection of nets will result in multiple sources");
                .
                .
                .

            }

            // Add all element connections of the new Net to this Net.
            net.nets.addAll(nnet.nets);
            net.elements.addAll(nnet.elements);
            net.ports.addAll(nnet.ports);
            net.porttypes.addAll(nnet.porttypes);
            net.inports += nnet.inports;
            net.outports += nnet.outports;
            net.tsoutports += nnet.tsoutports;

            // Add in the new ident list.
            net.idents.addAll(nnet.idents);

            // Append UCF and NCF strings from Net n to this Net.
            addConstraints("U", nnet.getConstraints("U"));
            addConstraints("N", nnet.getConstraints("N"));

            instances.remove(n.net);

            // Change all Nets associated with n.net to reference net.
            for (Net nn : nnet.nets)
                nn.net = net;
        }
    \end{lstlisting}
    This method connects two Nets. It does this by effectively
    merging the two NET subclasses so the the Nets share the same
    NET subclass; it adds the data from one NET to the other NET
    then deletes the former. It also traverses all Nets from the merged
    NET to change their nested NET class instance to the merged
    instance. If the argument Net is connected to GND or VCC then
    the method calls itself with the object and argument interchanged.
    This is to ensure that the Net identifier "GND" or "VCC"
    remains at the head of the list of connected identifiers and
    hence when {\bf setPrefIdent()} is called from TDEList.ncode()
    the "GND" or "VCC" will be encountered first in the list and will
    be set as the final preferred identifier (otherwise an "OPAD..."
    identifier might be encountered and prevail and the GND net in
    the netlist file would be labelled thus).             

\subsection{method addProperty()}

    \begin{lstlisting}[gobble=4, language=java]
        public void addProperty (String p, String v) {
            net.properties.put(p,  v);
        }
    \end{lstlisting}
    This method adds a property string to a net. P is the property identifier
    and v is the value string.

\subsection{Output Net instance to the EDIF netlist file}

    \begin{lstlisting}[gobble=4, language=java]
        static public void outputInstances (PrintWriter pw)
    \end{lstlisting}
    This method writes all Net instances to the EDIF file including
    all Element port connections.
